\begin{definition}
For a sample \(\omega\), vertex set \(I\), and tagged projected coordinate
\(e\), the predicate
\(\noderef{TwoBiteProjectionPairSameColorClosed}(\omega,I,e)\) is defined by
cases on the tag.  If \(e=\operatorname{red}(q)\), it means that there exists a
red base vertex \(r\) such that
\[
q\in
\noderef{TwoBiteBasePairs}
\bigl(\pi_R^\omega(\noderef{TwoBiteXPlus}(\omega,I,\operatorname{red}(r)))\bigr).
\]
If \(e=\operatorname{blue}(q)\), it means that there exists a blue base vertex
\(b\) such that
\[
q\in
\noderef{TwoBiteBasePairs}
\bigl(\pi_B^\omega(\noderef{TwoBiteXPlus}(\omega,I,\operatorname{blue}(b)))\bigr).
\]
There is no additional requirement in this predicate that \(e\) already belong
to \(\noderef{TwoBiteProjectionPairSet}(\omega,I)\).  It is the \(C^+(I)\)
part of the paper's closed-pair notation, recording same-color witnesses in
which the current projected pair is charged as the latest same-color edge.
\end{definition}
